\documentclass[12pt,a4paper]{article}
\usepackage[utf8]{inputenc}
\usepackage[T2A]{fontenc}
\usepackage[english,russian]{babel}
\usepackage{amsmath,amssymb,amsfonts}
\usepackage{geometry}
\usepackage{hyperref}
\geometry{margin=2.5cm}

\title{Angel-in-Pocket Shopping Assistant with GRA Nullification Engine\\
\large Bilingual Paper / Билингвальная статья}
\author{qqewq}
\date{\today}

\begin{document}
\maketitle

\begin{abstract}
\selectlanguage{english}
We present an AI-powered shopping assistant that detects risky products, suspicious reviews, and potentially hostile sellers by analyzing a local interaction graph and applying a GRA-inspired nullification engine. The system combines graph construction, recursive equilibrium computation, and a nullification cascade to produce product and seller trust scores. Subscription tiers control analysis depth and unlock richer diagnostics.

\selectlanguage{russian}
Представлен AI-ассистент для онлайн-шопинга, который выявляет рискованные товары, подозрительные отзывы и потенциально враждебных продавцов, анализируя локальный граф взаимодействий и применяя GRA-инспирированный движок обнуления. Система объединяет построение графа, рекурсивное вычисление равновесия и каскад обнуления для получения оценок доверия к товарам и продавцам. Уровни подписки управляют глубиной анализа и открывают более подробную диагностику.
\end{abstract}

\section{Introduction / Введение}

\selectlanguage{english}
Online marketplaces are increasingly affected by fake reviews, manipulative pricing, and adversarial behavior. A practical shopping assistant must therefore look beyond isolated product metadata and analyze the surrounding social and transactional context. In this paper, we describe a graph-based assistant that uses a GRA-style nullification engine to identify and dampen suspicious structures.

\selectlanguage{russian}
Онлайн-маркетплейсы всё сильнее страдают от фейковых отзывов, манипулятивного ценообразования и враждебного поведения. Практический шопинг-ассистент должен поэтому смотреть не только на отдельные характеристики товара, но и на окружающий социальный и транзакционный контекст. В этой статье мы описываем графовый ассистент, использующий GRA-движок обнуления для выявления и подавления подозрительных структур.

\section{System Overview / Обзор системы}

\selectlanguage{english}
The system is built around four components:
\begin{itemize}
    \item graph construction from products, sellers, reviews, and users;
    \item initial suspicion assignment using simple heuristics;
    \item recursive equilibrium propagation across the graph;
    \item nullification cascade to spread or attenuate risk.
\end{itemize}
The final output is a trust score in the interval $[0,1]$, where higher values indicate lower risk.

\selectlanguage{russian}
Система состоит из четырёх компонентов:
\begin{itemize}
    \item построение графа по товарам, продавцам, отзывам и пользователям;
    \item назначение начальной подозрительности на основе простых эвристик;
    \item рекурсивное распространение равновесия по графу;
    \item каскад обнуления для распространения или ослабления риска.
\end{itemize}
Итоговый результат — индекс доверия на интервале $[0,1]$, где большие значения означают меньший риск.
\end{document}